\documentclass{article}
\usepackage[default]{lora}
\PassOptionsToPackage{hyphens}{url}

\usepackage{polyglossia}
\setmainlanguage{english}
\setotherlanguage{russian}

\usepackage{microtype}
\usepackage{metalogo}
\usepackage[dvipsnames]{xcolor}

\usepackage[style=numeric, backend=biber]{biblatex}
\addbibresource{lora.bib}

\usepackage[
  colorlinks=true,
  linkcolor=MidnightBlue,
  citecolor=OliveGreen,
  urlcolor=NavyBlue,
  filecolor=Magenta
]{hyperref}

\usepackage{unicodefonttable}

\setcounter{secnumdepth}{0}

\makeatletter
\def\GetFileInfo#1{%
  \def\filename{#1}%
  \def\@tempb##1 ##2 ##3\relax##4\relax{%
    \def\filedate{##1}%
    \def\fileversion{##2}%
    \def\fileinfo{##3}}%
  \edef\@tempa{\csname ver@#1\endcsname}%
  \expandafter\@tempb\@tempa\relax? ? \relax\relax}
\makeatother

\GetFileInfo{lora.sty}

\newcommand{\abc}{abcdefghijklmnopqrstuvwxyz}
\newcommand{\ABC}{ABCDEFGHIJKLMNOPQRSTUVWXYZ}
\newcommand{\abv}{абвгдеёжзийклмнопрстуфхцчшщъыьэюя}
\newcommand{\ABV}{АБВГДЕЁЖЗИЙКЛМНОПРСТУФХЦЧШЩЪЫЬЭЮЯ}

% Alice's Adventures in Wonderland (1907)
\newcommand{\alicetext}{
  Alice was beginning to get very tired of sitting by her sister on the bank, and of having nothing to do: once or twice she had peeped into the book her sister was reading, but it had no pictures or conversations in it, ``and what is the use of a book,'' thought Alice, ``\textit{without pictures or conversations?}''
}

\begin{document}


\title{Sample of Lora fonts}
\author{Dmitriy Nosachev}
\date{Lora package version \fileversion, \filedate}
\maketitle
\tableofcontents

\clearpage

\section{Regular (weight = 400)}
\alicetext

\section{Bold (weight = 700)}
{\bfseries \alicetext}

\section{Medium (weight = 500)}
{\LoraMedium \alicetext}

\section{Semibold (weight = 600)}
{\LoraSemiBold \alicetext}

\section{Cyrillic}

\begin{russian}

% Письмо Н. В. Гоголю 15 июля 1847 г. (Белинский)
Не буду распространяться о вашем дифирамбе любовной связи русского народа с его владыками. Скажу прямо: этот дифирамб ни в ком не встретил себе сочувствия и уронил вас в глазах даже людей, в других отношениях очень близких к вам по их направлению. Что касается до меня лично, предоставляю вашей совести упиваться созерцанием божественной красоты самодержавия (оно покойно, да — и выгодно), только продолжайте благоразумно созерцать его из вашего \textit{прекрас\-ного далека}: вблизи-то оно не так прекрасно и не так безопасно…

\end{russian}


\section{Alphabets}


{\setlength{\parindent}{0pt}
\setlength{\parskip}{1ex}


\ABC\\
\textit{\ABC}\\
\textbf{\ABC}\\
\textbf{\itshape \ABC}\\
\ABV\\
\textit{\ABV}\\


\abc\\
\textit{\abc}\\
\textbf{\abc}\\
\textbf{\itshape \abc}\\
\abv\\
\textit{\abv}\\
\textbf{\abv}\\
\textbf{\itshape \abv}\\
0123456789}

\section{Font table}


\displayfonttable[color=BrickRed]{Lora[wght].ttf}

\nocite{*}

\clearpage
\printbibliography


\end{document}
